\section{An intrinsic capacity characterization of the strong first-energy class}
\label{sec:r11-intrinsic-first-energy-characterization}

Retain the strict analytic box, original background extension and
arbitrary original $S$ of (TC1)--(TC2). Let
$\mathcal A_S^{[-R,0]}(V)$ denote the global bounded capacitary
completion of Lemma~\ref{lem:r11-analytic-bounded-completion}:
its members are Borel limits, in every local outer capacity,
of actual bounded classical functions on the same whole domain
$V$, with values in $[-R,0]$.
The word ``global'' specifies that admissibility domain; it
does not assert a uniform convergence on all of $V$.

The theorem below replaces a strongly Cauchy approximating row
by an intrinsic condition on all bounded stages and one
capacity integral. The bounded-stage admissibility condition
is retained explicitly. An arbitrary Borel function with a
finite numerical capacity integral is not asserted admissible.

\begin{lemma}[Ordered discrepancy is bounded by energy loss]
\label{lem:r11-ordered-completed-discrepancy}
If $a,b\in\widehat{\mathcal E}^1_{\rm str}(V)$ and
$a\leq b$ outside a locally outer-capacity-null set, then
\[
       0\leq\widehat I(a,b)
                            \leq\widehat e(a)-\widehat e(b).
 \tag{CI1}
\]
\end{lemma}

\noindent\textbf{Status: \texttt{DIRECT\_COROLLARY}} of the
proved completed symmetric mixed energy and its positivity.

\begin{proof}
The completed mixed energy is decreasing in each potential
with respect to pointwise order: in the undifferentiated
slot this is positivity of the product measure, and symmetry
gives it in every slot. All products ignore the exceptional
outer-null set. Therefore
\[
       \widehat B(a,b^{[q]})\geq\widehat e(b),\qquad
       \widehat B(b,a^{[q]})\geq\widehat e(b).
\]
Insert these two inequalities in the definition (IC17) of
$\widehat I$. Its nonnegativity is already proved. This gives
(CI1), with finite global integrals throughout.
\end{proof}

\begin{theorem}[Bounded-stage admissibility and one capacity integral]
\label{thm:r11-intrinsic-strong-first-energy}
For an actual Borel function $u:V\to[-\infty,0]$, the following
are equivalent:
\begin{enumerate}
\item $u\in\widehat{\mathcal E}^1_{\rm str}(V)$.
\item For every positive integer $m$,
$u^{(m)}=\max(u,-m)$ belongs to
$\mathcal A_S^{[-m,0]}(V)$, and
\[
                J_1^*(u):=\int_0^\infty
                       t^q C_S^*(\{u<-t\};V)\,dt<\infty.
 \tag{CI2}
\]
\end{enumerate}
In the second condition no current of the unbounded function
is assumed in advance. Once either condition holds, the
canonical original currents and all critical first-energy
identities are those of (IC10)--(IC20), and
\[
       2^{-q-1}J_1^*(u)\leq\widehat e(u)
                                  \leq2^{q+1}J_1^*(u).
 \tag{CI3}
\]
The capacity integral is unchanged if raw capacity is used.
This last equality is a conclusion for the resulting completed
strict sublevels, not a premise of the second condition.
\end{theorem}

\noindent\textbf{Status: \texttt{NEW\_PROVED}} as a converse
capacity characterization of the full strong Borel class,
with its bounded-stage admissibility gate retained.

\begin{proof}
Suppose first that $u$ belongs to the strong completion.
An actual strongly Cauchy defining row, clipped at $-m$,
is a sequence of actual global bounded classical functions
in $[-m,0]$ converging to $u^{(m)}$ in the required local
outer capacities. This proves every bounded-stage condition.
The completed capacity-layer theorem (CY3) at $s=1$ gives
(CI2) and the lower inequality of (CI3).

Conversely assume precisely the second condition. Let $\rho$
be the actual barrier (TC3), so $-1\leq\rho<0$ on $V$ and
$\rho$ has finite total top mass. For each integer $j\geq1$
define the Borel function
\[
                     v_j=\max(u,j\rho)
                         =\max(u^{(j)},j\rho).
 \tag{CI4}
\]
The equality follows from $j\rho\geq-j$.
Choose global actual bounded defining rows for $u^{(j)}$.
Taking their maximum with $j\rho$ gives actual classical
functions between $j\rho$ and zero, converging to $v_j$
in the bounded local capacity topology. Each such function
belongs to $\mathcal E^1$ by upward stability above the
actual first-energy barrier $j\rho$. Thus
$v_j\in\widehat{\mathcal E}^1_{j\rho}$ by the definition
of the barrier completion, and hence belongs to the strong
completion by its proved embedding. This step defines
products only for $v_j$, not yet for the given $u$.

Since $v_j\geq u$, monotonicity of outer capacity gives
$\widehat J_1(v_j)\leq J_1^*(u)$.
The already proved theorem (CY3), applied to the valid
completed member $v_j$, therefore gives
\[
                    \widehat e(v_j)\leq2^{q+1}J_1^*(u).
 \tag{CI5}
\]
Pointwise $v_{j+1}\leq v_j$. Hence their energies increase
and are bounded by (CI5). For $l\geq j$, (CI1) gives
\[
                  \widehat I(v_l,v_j)
                           \leq\widehat e(v_l)-\widehat e(v_j).
 \tag{CI6}
\]
The increasing bounded energy sequence has a finite limit,
so $(v_j)$ is strongly Cauchy. Completeness (IC20) supplies
a strong completed limit $v$.

It remains to identify its actual capacity class with $u$;
pointwise convergence alone is not used for that purpose.
Fix a compact $H\Subset V$ and a clipping depth $R>0$.
The continuous strict negativity of $\rho$ makes
$\max_H\rho<0$. For all sufficiently large $j$, one has
$j\rho<-R$ on $H$, and therefore the exact scalar identity
\[
                              v_j^{(R)}=u^{(R)}
                                      \quad\hbox{on }H
 \tag{CI7}
\]
holds. Strong convergence gives local clipped capacity
convergence to $v$. Equation (CI7) identifies that limit
with $u^{(R)}$ on every such compact core. The countable
core and integer-depth definition of the completion thus
identifies $u$ with the strong completed limit, using its
specified Borel representative. This proves membership.

The full completed current and energy calculus now applies
to $u$. Applying (CY3) at $s=1$ gives both inequalities in
(CI3); the strict-gap regularity preceding (CC8) gives the
raw/outer equality for its sublevels. No unbounded current
of $u$ was used before it was constructed through completion.
\end{proof}

\begin{corollary}[Monotone closure with the exact energy hypothesis]
\label{cor:r11-strong-monotone-closure}
Let $u_j$ be members of
$\widehat{\mathcal E}^1_{\rm str}(V)$, monotone outside one
locally outer-null set, and let $u$ be their actual pointwise
monotone limit there, with any nonpositive extended Borel values
on that null set.
If $u_j$ increases, then $u$ belongs to the strong completion
and $\widehat I(u_j,u)\to0$.
If $u_j$ decreases, the same conclusion holds provided
\[
                            \sup_j\widehat e(u_j)<\infty.
 \tag{CI8}
\]
Every original critical local current and global mixed
weighted-measure convergence assertion (IC11)--(IC13)
then holds for finite tuples satisfying these hypotheses.
Ordinary upper semicontinuous regularization of the limit
is not part of this conclusion.
\end{corollary}

\noindent\textbf{Status: \texttt{NEW\_PROVED}} for monotone
sequences in the full strong completion; the increasing
bounded relative precursor was already frozen in R2.

\begin{proof}
For an increasing sequence, the nonnegative energies decrease
and are bounded above by the first energy. For a decreasing
sequence they increase, with the explicit bound (CI8).
In either case they converge. Equation (CI1), applied in
the appropriate order to two terms, shows that the sequence
is strongly Cauchy. Completeness gives a strong limit $v$.

Choose a subsequence whose fixed bounded truncations converge
pointwise off a locally outer-null set, using the usual
summable-capacity diagonal construction over compact cores
and integer depths. Monotonicity makes that pointwise limit
equal to the specified $u$. Thus $u=v$ in the completed
capacity class, and the entire sequence converges strongly.
The strong continuity theorem gives all the stated current
and weighted-measure limits. No current equality outside its
original test sheaf is asserted.
\end{proof}

\paragraph{Exact scope.}
The intrinsic condition (CI2) still requires every bounded
truncation to be an admissible global capacitary limit of
actual bounded classical functions. It therefore does not
turn a numerical capacity bound into ambient psh geometry.
It removes a prescribed actual lower barrier and a prescribed
strong approximating row from the characterization of the
completed first-energy class, while retaining the original
background, absolute value and all original tests. It does
not prove arbitrary-source Dirichlet existence or general
raw-capacity outer regularity.
